% The second eigenvalue of random regular graphs.
% Compile from this folder with:  latexmk -pdf slides.tex   (or pdflatex slides.tex, twice)
\documentclass[aspectratio=169,10pt]{beamer}

\usepackage[T1]{fontenc}
\usepackage[utf8]{inputenc}
\usepackage{lmodern}
\usepackage{amsmath,amssymb,mathtools}
\usepackage{booktabs,array}
\newcolumntype{L}[1]{>{\raggedright\arraybackslash}p{#1}}
\usepackage{pgfplots}
\usepgfplotslibrary{groupplots}
\pgfplotsset{compat=1.16}

% ---------- appearance ----------
\usetheme{default}
\usefonttheme[onlymath]{serif}
\definecolor{ink}{HTML}{1B2A41}
\definecolor{accent}{HTML}{A23B1E}
\definecolor{barfill}{HTML}{C9D3DF}
\setbeamercolor{structure}{fg=ink}
\setbeamercolor{normal text}{fg=black!88}
\setbeamercolor{frametitle}{fg=ink}
\setbeamercolor{alerted text}{fg=accent}
\setbeamercolor{block title}{fg=ink, bg=black!7}
\setbeamercolor{block body}{bg=black!3}
\setbeamerfont{frametitle}{size=\Large, series=\bfseries}
\setbeamerfont{title}{size=\huge, series=\bfseries}
\setbeamerfont{subtitle}{size=\large}
\setbeamerfont{block title}{size=\normalsize, series=\bfseries}
\setbeamerfont{block body}{size=\normalsize}
\setbeamertemplate{navigation symbols}{}
\setbeamertemplate{itemize items}[circle]
\setbeamertemplate{enumerate items}[default]
\setbeamertemplate{frametitle}{\vspace{5mm}\insertframetitle\par\vspace{2mm}}
\setbeamertemplate{footline}{\hfill{\color{black!45}\footnotesize\insertframenumber}\hspace{7mm}\vspace{4mm}}
\setbeamersize{text margin left=14mm, text margin right=14mm}
\setlength{\parskip}{0.6em}
% Extra air around definition and theorem blocks.
\addtobeamertemplate{block begin}{\vspace{1.2mm}}{\vspace{0.4mm}}
\addtobeamertemplate{block end}{\vspace{0.4mm}}{\vspace{1.2mm}}

% Statements are upright; only the label is emphasised.
\setbeamertemplate{theorems}[default]
\setbeamerfont{block body}{shape=\upshape}
\newtheorem{proposition}[theorem]{Proposition}

% ---------- figure sizes ----------
\newcommand{\specw}{4.8cm}
\newcommand{\spech}{2.3cm}
\newcommand{\specsep}{0.5cm}
\newcommand{\nbw}{5.5cm}
\newcommand{\dyckunit}{3.9mm}
\newcommand{\shapeunit}{5.2mm}
\newcommand{\shapesep}{4.2cm}
\newcommand{\tangleunit}{4.6mm}
\newcommand{\dumbunit}{3.5mm}
\newcommand{\liftunit}{7mm}

% ---------- notation ----------
\DeclareMathOperator{\tr}{tr}
\newcommand{\edge}{2\sqrt{d-1}}
\newcommand{\eps}{\varepsilon}
\newcommand{\one}{\mathbf{1}}
\newcommand{\E}{\mathbb{E}}
\newcommand{\intuition}{{\color{accent}\textbf{Intuition.}}\ }
\newcommand{\pf}{\textbf{Proof sketch.}\ }

\title{The second eigenvalue of random regular graphs}
\author{}
\date{}

\begin{document}

\begin{frame}[plain]
  \vfill
  {\usebeamerfont{title}\usebeamercolor[fg]{title} The second eigenvalue of\\[1mm] random regular graphs\par}
  \vspace{8mm}
  {\color{black!60} The Alon--Boppana bound, Friedman's theorem and the trace method\par}
  \vfill
\end{frame}

% ------------------------------------------------------------------
\begin{frame}{Outline}
  \begin{enumerate}
    \setlength{\itemsep}{2.5mm}
    \item \textbf{Spectra of regular graphs.} Adjacency spectrum, spectral gap, Cheeger's inequality.
    \item \textbf{The barrier $\edge$.} Closed walks in the tree; the Alon--Boppana bound.
    \item \textbf{Random regular graphs.} The Kesten--McKay law; Friedman's theorem.
    \item \textbf{The obstruction.} The trace method, subgraph counts, tangles.
    \item \textbf{The resolution.} The non-backtracking matrix; tangle-free graphs.
    \item \textbf{Beyond the literature.} Graph lifts; recovering a hidden partition.
  \end{enumerate}
\end{frame}

% ==================================================================
\section{Spectra of regular graphs}

\begin{frame}{Adjacency spectrum}
  \begin{definition}
    Let $G$ be a $d$-regular graph on $n$ vertices with adjacency matrix $A$ and eigenvalues
    $d=\lambda_1\ge\lambda_2\ge\dots\ge\lambda_n\ge-d$. The eigenvalue $\lambda_1=d$, with
    eigenvector $\one$, is called \emph{trivial}. Write
    $\lambda(G)=\max\bigl(\lambda_2,\,|\lambda_n|\bigr)$.
  \end{definition}
  \begin{proposition}[Walk counting]
    $(A^k)_{uv}$ is the number of walks of length $k$ from $u$ to $v$. Consequently
    \[
      \tr(A^k)\;=\;\sum_{i=1}^{n}\lambda_i^{k}\;=\;\#\{\text{closed walks of length }k\}.
    \]
    \vspace{-5mm}
  \end{proposition}
  \textbf{Example.} Triangle: spectrum $\{2,-1,-1\}$;\; $2^3-1-1=6$ closed walks of length $3$.

  \intuition Questions about eigenvalues become questions about counting walks.
\end{frame}

% ------------------------------------------------------------------
\begin{frame}{Spectral gap and edge expansion}
  \begin{definition}
    The \emph{spectral gap} of $G$ is $d-\lambda_2$. The \emph{edge expansion} of $G$ is
    $h(G)=\min\limits_{0<|S|\le n/2} e(S,\bar S)/|S|$, where $e(S,\bar S)$ counts the edges leaving $S$.
    A sequence of $d$-regular graphs is an \emph{expander family} if $h\ge c>0$ along the sequence.
  \end{definition}
  \begin{proposition}[Dirichlet form]
    For every $x\in\mathbb{R}^{V}$,\quad
    $d\,\lVert x\rVert^2-x^{\mathsf T}Ax=\sum_{\{u,v\}\in E}(x_u-x_v)^2$.\quad Hence
    \[
      d-\lambda_2\;=\;\min_{x\perp\one}\ \frac{\sum_{\{u,v\}\in E}(x_u-x_v)^2}{\sum_u x_u^2}.
    \]
    \vspace{-5mm}
  \end{proposition}
  \intuition The gap is small exactly when some mean-zero function is nearly constant across every edge.
\end{frame}

% ------------------------------------------------------------------
\begin{frame}{Cheeger's inequality}
  \begin{theorem}[Cheeger; Dodziuk, Alon--Milman]
    For every $d$-regular graph $G$,\qquad
    $\dfrac{h(G)^2}{2d}\;\le\;d-\lambda_2\;\le\;2\,h(G)$.
  \end{theorem}
  \begin{columns}[T,onlytextwidth]
    \begin{column}{0.50\textwidth}
      \textbf{Examples.}\\[1mm]
      {\small
      \begin{tabular}{@{}lll@{}}
        \toprule
        Graph on $n$ vertices & Edges & $h$\\
        \midrule
        Complete graph & $\approx n^2/2$ & $\approx n/2$\\
        Cycle & $n$ & $\approx 4/n$\\
        Complete binary tree & $n-1$ & $\approx 2/n$\\
        Expander family & $dn/2$ & $\ge c$\\
        \bottomrule
      \end{tabular}}
    \end{column}
    \begin{column}{0.40\textwidth}
      \centering
      % Two dense clusters joined by a single edge.
% Requires: \dumbunit (unit length), colours ink, accent and barfill.
\begin{tikzpicture}[x=\dumbunit, y=\dumbunit]
  \draw[accent, thick] (-1.1,0) -- (1.1,0);
  \foreach \s/\lab in {-1/{+1}, 1/{-1}} {
    \fill[barfill] (2.1*\s,0) circle (1.05);
    \draw[ink] (2.1*\s,0) circle (1.05);
    \node[font=\footnotesize, anchor=north] at (2.1*\s,-1.15) {$x=\lab$};
  }
  \fill[ink] (-1.1,0) circle (1.4pt) (1.1,0) circle (1.4pt);
\end{tikzpicture}


      \raggedright\small
      One edge between two clusters: the test vector $x=\pm1$ gives $d-\lambda_2\le4/n$.
    \end{column}
  \end{columns}
  \vspace{2mm}
  \intuition A spectral gap certifies the absence of bottlenecks. How large can it be?
\end{frame}

% ==================================================================
\section{The barrier}

\begin{frame}{Closed walks in the infinite tree}
  \begin{lemma}
    Let $t_{2k}$ be the number of closed walks of length $2k$ from a fixed vertex of the infinite
    $d$-regular tree $T_d$. Then $t_{2k}\ge C_k\,(d-1)^k$, where $C_k\sim 4^k/(\sqrt{\pi}\,k^{3/2})$ is the
    Catalan number, and $\lim_{k\to\infty}t_{2k}^{1/2k}=\edge$.
  \end{lemma}
  \begin{columns}[T,onlytextwidth]
    \begin{column}{0.42\textwidth}
      \centering
      % Distance from the starting vertex along a closed walk of length 12 in a tree.
% Requires: \dyckunit (length of one step), colours ink and accent.
\begin{tikzpicture}[x=\dyckunit, y=\dyckunit, line cap=round, line join=round]
  \draw[black!35] (0,0) -- (12,0);
  \foreach \x/\y/\xx/\yy in {0/0/1/1, 1/1/2/2, 3/1/4/2, 4/2/5/3, 8/0/9/1, 9/1/10/2}
    \draw[accent, very thick] (\x,\y) -- (\xx,\yy);
  \foreach \x/\y/\xx/\yy in {2/2/3/1, 5/3/6/2, 6/2/7/1, 7/1/8/0, 10/2/11/1, 11/1/12/0}
    \draw[ink, very thick] (\x,\y) -- (\xx,\yy);
  \foreach \x/\y in {0/0,1/1,2/2,3/1,4/2,5/3,6/2,7/1,8/0,9/1,10/2,11/1,12/0}
    \fill[black!70] (\x,\y) circle (1.1pt);
  \node[font=\footnotesize, anchor=north] at (6,-0.25) {step};
  \node[font=\footnotesize, rotate=90, anchor=south] at (-0.4,1.5) {distance};
\end{tikzpicture}

    \end{column}
    \begin{column}{0.54\textwidth}
      \pf A closed walk takes $k$ steps \textcolor{accent}{away} from its start and $k$ steps
      \textcolor{ink}{back}. The pattern of the two is a Dyck path ($C_k$ choices); each step away
      has at least $d-1$ continuations; each step back is forced.
    \end{column}
  \end{columns}
  \vspace{2mm}
  \intuition In $\edge$, the factor $2$ counts \emph{when} to turn around and $\sqrt{d-1}$ counts \emph{where} to go.
\end{frame}

% ------------------------------------------------------------------
\begin{frame}{The Alon--Boppana bound}
  \begin{theorem}[Alon--Boppana]
    Every $d$-regular graph $G$ on $n$ vertices satisfies\quad $\lambda_2(G)\ge\edge-o_n(1)$.
  \end{theorem}
  \pf (for $\lambda(G)$)
  \begin{enumerate}
    \item $T_d$ covers $G$, so closed walks in $T_d$ project injectively to closed walks in $G$.
    \item Hence\quad $d^{2k}+(n-1)\,\lambda(G)^{2k}\;\ge\;\tr(A^{2k})\;\ge\;n\,C_k\,(d-1)^k$.
    \item Take $k=c\log n$ with $c$ small, so that $d^{2k}\ll n\,(\edge)^{2k}$; then take $2k$-th roots.
  \end{enumerate}
  \begin{definition}
    A $d$-regular graph $G$ is \emph{Ramanujan} if $\lambda(G)\le\edge$.
  \end{definition}
  \intuition The tree is extremal: every graph has at least as many closed walks.
\end{frame}

% ==================================================================
\section{Random regular graphs}

\begin{frame}{Random regular graphs}
  \begin{definition}
    $G_{n,d}$ denotes a uniformly random simple $d$-regular graph on $n$ vertices ($nd$ even).
  \end{definition}
  \begin{theorem}[McKay, 1981]
    For fixed $d\ge2$, the empirical eigenvalue distribution of $G_{n,d}$ converges to the
    Kesten--McKay law, the spectral measure of $T_d$, which is supported on $[-\edge,\,\edge]$.
  \end{theorem}
  \vspace{-1mm}
  \centering
  % Eigenvalue histograms of random d-regular graphs against the Kesten--McKay density.
% Requires: \specw, \spech (panel size), \specsep (gap), colours barfill and accent.
% The horizontal variable is t = lambda / (2 sqrt(d-1)), so the bulk is [-1, 1] in every panel.
\begin{tikzpicture}
\begin{groupplot}[
  group style={group size=3 by 1, horizontal sep=\specsep},
  width=\specw, height=\spech,
  xmin=-1.18, xmax=1.18, ymin=0, ymax=0.80,
  axis x line*=bottom, axis y line=none,
  xtick={-1,1}, tick align=outside, tick style={black!60},
  x axis line style={black!60},
  title style={font=\small, yshift=-1ex},
  ticklabel style={font=\footnotesize},
  clip=false,
]
\nextgroupplot[title={$d=3$}, xticklabels={$-2\sqrt{2}$, $2\sqrt{2}$}]
  \addplot[ybar, bar width=0.05, draw=white, line width=0.2pt, fill=barfill] table {data/hist_d3.dat};
  \addplot[accent, thick, domain=-1:1, samples=241]
    {2.828427*3*sqrt(max(0, 8 - (2.828427*x)^2)) / (2*pi*(9 - (2.828427*x)^2))};
\nextgroupplot[title={$d=6$}, xticklabels={$-2\sqrt{5}$, $2\sqrt{5}$}]
  \addplot[ybar, bar width=0.05, draw=white, line width=0.2pt, fill=barfill] table {data/hist_d6.dat};
  \addplot[accent, thick, domain=-1:1, samples=241]
    {4.472136*6*sqrt(max(0, 20 - (4.472136*x)^2)) / (2*pi*(36 - (4.472136*x)^2))};
\nextgroupplot[title={$d=16$}, xticklabels={$-2\sqrt{15}$, $2\sqrt{15}$}]
  \addplot[ybar, bar width=0.05, draw=white, line width=0.2pt, fill=barfill] table {data/hist_d16.dat};
  \addplot[accent, thick, domain=-1:1, samples=241]
    {7.745967*16*sqrt(max(0, 60 - (7.745967*x)^2)) / (2*pi*(256 - (7.745967*x)^2))};
\end{groupplot}
\end{tikzpicture}


  \vspace{-1mm}
  {\scriptsize\color{black!60} One sample with $n=2000$ per degree (bars) against the limiting density (curve).}
\end{frame}

% ------------------------------------------------------------------
\begin{frame}{Friedman's theorem}
  \begin{theorem}[Friedman, 2003; conjectured by Alon, 1986]
    For every $d\ge3$ and $\eps>0$,
    \[
      \Pr\bigl[\,\lambda(G_{n,d})\;\ge\;\edge+\eps\,\bigr]\;\longrightarrow\;0
      \qquad(n\to\infty).
    \]
    \vspace{-5mm}
  \end{theorem}
  \vspace{1mm}
  \textbf{Remarks.}
  \begin{itemize}
    \item By Alon--Boppana the constant $\edge$ is optimal: almost every $d$-regular graph is almost Ramanujan.
    \item McKay's theorem controls the bulk of the spectrum, but allows $o(n)$ eigenvalues anywhere.
          Friedman's theorem says that \alert{none} lies outside the bulk.
  \end{itemize}
  \intuition Locally $G_{n,d}$ looks like $T_d$, so it should have about as few closed walks as a graph can.
\end{frame}

% ==================================================================
\section{The obstruction}

\begin{frame}{The trace method}
  \begin{lemma}
    Let $\rho=\edge$ and $k\ge1$. If\/ $\sum_{i\ge2}\lambda_i^{2k}\le n\,(\rho+\eps)^{2k}$, then\/
    $\lambda(G)\le n^{1/2k}\,(\rho+\eps)$.
  \end{lemma}
  \pf $\lambda(G)^{2k}$ is one of the non-negative terms of the sum.\hfill$\square$

  \vspace{1mm}
  The bound is useful only if $n^{1/2k}\to1$, that is, $k\gg\log n$.

  \vspace{1mm}
  \begin{center}\small
  \begin{tabular}{@{}lll@{}}
    \toprule
    & Walk length & What the walks see\\
    \midrule
    Lower bound (Alon--Boppana) & $k\asymp\log n$ & only the tree\\
    Upper bound (Friedman) & $k\gg\log n$ & the whole graph, many times over\\
    \bottomrule
  \end{tabular}
  \end{center}
  \intuition A high power is an amplifier: it lets one eigenvalue outweigh $n$ others.
\end{frame}

% ------------------------------------------------------------------
\begin{frame}{Subgraph counts}
  \begin{lemma}
    Let $H$ be a fixed connected graph with $v$ vertices, $e$ edges and maximum degree at most $d$.
    The expected number of copies of $H$ in a random $d$-regular graph on $n$ vertices is $\Theta(n^{\,v-e})$.
  \end{lemma}
  \vspace{-1mm}
  \begin{center}
    % Three shapes a closed walk can trace: a tree, a cycle, and two cycles sharing a vertex.
% Requires: \shapeunit (unit length), \shapesep (gap between shapes), colours ink and accent.
\begin{tikzpicture}[x=\shapeunit, y=\shapeunit,
    v/.style={circle, fill=ink, inner sep=1.6pt}, e/.style={ink, semithick}]
  % tree: v = 5, e = 4
  \begin{scope}[yshift={-1*\shapeunit}]
    \node[v] (a) at (0,2) {}; \node[v] (b) at (-1,1) {}; \node[v] (c) at (1,1) {};
    \node[v] (d) at (-1.5,0) {}; \node[v] (f) at (-0.5,0) {};
    \draw[e] (a)--(b) (a)--(c) (b)--(d) (b)--(f);
    \node[font=\footnotesize, align=center, anchor=north] at (0,-0.45) {tree\\ $v-e=1$\\ $\asymp n$ copies};
  \end{scope}
  % cycle: v = 5, e = 5
  \begin{scope}[xshift=\shapesep]
    \foreach \i in {0,...,4} \node[v] (p\i) at ({90+72*\i}:1.05) {};
    \draw[e] (p0)--(p1)--(p2)--(p3)--(p4)--(p0);
    \node[font=\footnotesize, align=center, anchor=north] at (0,-1.45) {one cycle\\ $v-e=0$\\ $\asymp 1$ copies};
  \end{scope}
  % two triangles sharing a vertex: v = 5, e = 6
  \begin{scope}[xshift={2*\shapesep}]
    \node[v, fill=accent] (o) at (0,0) {};
    \node[v] (l1) at (-1.3,0.75) {}; \node[v] (l2) at (-1.3,-0.75) {};
    \node[v] (r1) at (1.3,0.75) {};  \node[v] (r2) at (1.3,-0.75) {};
    \draw[e] (o)--(l1)--(l2)--(o) (o)--(r1)--(r2)--(o);
    \node[font=\footnotesize, align=center, anchor=north] at (0,-1.45) {two cycles\\ $v-e=-1$\\ $\asymp 1/n$ copies};
  \end{scope}
\end{tikzpicture}

  \end{center}
  \vspace{-2mm}
  \intuition Each independent cycle costs a factor of $n$. Closed walks that trace trees dominate
  $\E\,\tr(A^{2k})$, which suggests $\E\sum_{i\ge2}\lambda_i^{2k}\approx n\,(\edge)^{2k}$.
\end{frame}

% ------------------------------------------------------------------
\begin{frame}{Tangles}
  \begin{definition}
    For a finite connected graph $\psi$, let $h(\psi)$ be the exponential growth rate of the number
    of non-backtracking walks of length $k$ in $\psi$. Call $\psi$ a \emph{tangle} (for degree $d$) if
    $h(\psi)>\sqrt{d-1}$.
  \end{definition}
  \begin{theorem}[Friedman]
    Let $\psi$ be a tangle and $h=h(\psi)$. Every $d$-regular graph $G$ on $n$ vertices that contains $\psi$ satisfies
    \[
      \lambda_2(G)\;\ge\;h+\frac{d-1}{h}-o_n(1)\;>\;\edge .
    \]
    \vspace{-5mm}
  \end{theorem}
  \textbf{Example.} A vertex with $m$ self-loops has $h=2m-1$. It is a tangle when $2m-1>\sqrt{d-1}$,
  and occurs in the permutation model with probability $\asymp n^{1-m}$.
\end{frame}

% ------------------------------------------------------------------
\begin{frame}{Tangles: proof sketch and a numerical check}
  \begin{columns}[T,onlytextwidth]
    \begin{column}{0.15\textwidth}
      \centering\vspace{1mm}
      % A vertex with m = 3 self-loops in a d-regular graph, with trees hanging below it.
% Requires: \tangleunit (unit length), colours ink and accent.
\begin{tikzpicture}[x=\tangleunit, y=\tangleunit,
    v/.style={circle, fill=ink, inner sep=1.3pt}, e/.style={ink, semithick}]
  % Explicit bounding box: the loops' control points would otherwise inflate it.
  \path[use as bounding box] (-1.7,-2.7) rectangle (1.7,1.5);
  \node[circle, fill=accent, inner sep=2pt] (o) at (0,0) {};
  \foreach \a in {40,90,140}
    \draw[accent, thick] (o) to[out=\a-24, in=\a+24, looseness=38] (o);
  \node[v] (a) at (-0.9,-1.1) {}; \node[v] (b) at (0.9,-1.1) {};
  \node[v] (a1) at (-1.35,-2.1) {}; \node[v] (a2) at (-0.45,-2.1) {};
  \node[v] (b1) at (0.45,-2.1) {};  \node[v] (b2) at (1.35,-2.1) {};
  \draw[e] (o)--(a) (o)--(b) (a)--(a1) (a)--(a2) (b)--(b1) (b)--(b2);
  \foreach \p in {a1,a2,b1,b2} \draw[e, densely dotted] (\p) -- ++(0,-0.5);
\end{tikzpicture}

    \end{column}
    \begin{column}{0.47\textwidth}
      \pf (vertex $o$ with $m$ loops)
      \begin{enumerate}
        \item Let $f(v)=h^{-\operatorname{dist}(v,o)}$ within radius $r$ of $o$, and $0$ beyond.
              Then $Af\ge\bigl(h+\tfrac{d-1}{h}\bigr)f$ inside the ball.
        \item $\lVert f\rVert$ stays bounded as $r\to\infty$ because $(d-1)/h^2<1$.
        \item $f$ lives on $O(1)$ vertices; comparing with the indicator of the far vertices gives the bound on $\lambda_2$.
      \end{enumerate}
    \end{column}
    \begin{column}{0.30\textwidth}
      \begin{block}{Numerical check}
        \small
        $d=8$,\; $n=2000$,\; $m=3$\\[3pt]
        \begin{tabular}{@{}l@{\hspace{3mm}}r@{}}
          $2\sqrt7$ & $5.292$\\
          $\lambda_2$, random graph & $5.251$\\
          $\lambda_2$, loops planted & $6.397$\\
          $h+(d-1)/h$ & $6.400$\\
        \end{tabular}
      \end{block}
    \end{column}
  \end{columns}
  \vspace{2mm}
  \intuition A tangle recirculates walks faster than the surrounding trees can carry them away.
\end{frame}

% ------------------------------------------------------------------
\begin{frame}{Failure of the expected trace}
  \begin{corollary}
    Suppose a tangle $\psi$ occurs in the random graph $G$ with probability at least $a\,n^{-c}$, and let
    $\lambda_\psi=h+\frac{d-1}{h}>\rho=\edge$. Then, for small $\eps>0$ and large $n$:
    \begin{enumerate}
      \item $\Pr\bigl[\lambda_2(G)>\rho+\eps\bigr]\;\ge\;a\,n^{-c}$;
      \item $\E\sum_{i\ge2}\lambda_i^{2k}\;\ge\;a\,n^{-c}\,(\lambda_\psi-\eps)^{2k}$, which exceeds
            $n\,(\rho+\eps)^{2k}$ as soon as $k\ge C\log n$.
    \end{enumerate}
  \end{corollary}
  \vspace{1mm}
  \textbf{Consequences.}
  \begin{itemize}
    \item The convergence in Friedman's theorem is at best polynomial in $n$.
    \item No bound on $\E\,\tr(A^{2k})$ with $k\gg\log n$ can prove the theorem.
  \end{itemize}
  \intuition Tangles are rare, but the amplifier that isolates $\lambda_2$ amplifies them too.
\end{frame}

% ==================================================================
\section{The resolution}

\begin{frame}{The non-backtracking matrix}
  \begin{columns}[c,onlytextwidth]
    \begin{column}{0.40\textwidth}
      \centering
      % Spectrum of the non-backtracking matrix of a random 4-regular graph on 150 vertices.
% Requires: \nbw (plot width), colours ink and accent.
\begin{tikzpicture}
\begin{axis}[
  width=\nbw, axis equal image,
  xmin=-2.3, xmax=3.6, ymin=-2.15, ymax=2.15,
  axis lines=middle, axis line style={black!45},
  xtick={-1,1,3}, ytick=\empty,
  ticklabel style={font=\footnotesize, fill=white, inner sep=1pt},
  tick style={black!45},
  xlabel={$\operatorname{Re}\mu$}, ylabel={$\operatorname{Im}\mu$},
  xlabel style={font=\footnotesize, at={(ticklabel* cs:1)}, anchor=west},
  ylabel style={font=\footnotesize, at={(ticklabel* cs:1)}, anchor=south},
  clip=false,
]
  \draw[accent, dashed] (axis cs:0,0) circle [radius=1.7320508];
  \addplot[only marks, mark=*, mark size=0.9pt, ink] table {data/nb_d4_n150.dat};
  \node[font=\footnotesize, accent, anchor=south west] at (axis cs:1.25,1.25) {$|\mu|=\sqrt{d-1}$};
  \node[font=\footnotesize, anchor=south] at (axis cs:3,0.15) {trivial};
\end{axis}
\end{tikzpicture}


      {\scriptsize\color{black!60} Spectrum of $B$ for a random $4$-regular graph, $n=150$.}
    \end{column}
    \begin{column}{0.57\textwidth}
      \begin{definition}
        $B$ is the matrix on the $nd$ directed edges of $G$ with $B_{(u\to v),(v\to w)}=1$ if $w\neq u$, and $0$ otherwise.
      \end{definition}
      \begin{theorem}[Ihara--Bass, $d$-regular case]
        The eigenvalues of $B$ are $\pm1$ and, for each eigenvalue $\lambda$ of $A$, the two roots of
        $\mu^2-\lambda\mu+(d-1)=0$.
      \end{theorem}
      \textbf{Corollary.} $G$ is Ramanujan iff every non-trivial eigenvalue of $B$ has $|\mu|=\sqrt{d-1}$.
    \end{column}
  \end{columns}
\end{frame}

% ------------------------------------------------------------------
\begin{frame}{Tangle-free graphs}
  \begin{definition}[Bordenave]
    $G$ is \emph{$\ell$-tangle-free} if every ball of radius $\ell$ in $G$ contains at most one cycle.
  \end{definition}
  \textbf{Lemma.} A random $d$-regular graph is $\ell$-tangle-free with probability $1-O\bigl((d-1)^{4\ell}/n\bigr)$.
  \begin{proposition}[Strategy]
    Let $B^{(\ell)}$ count only the tangle-free non-backtracking walks of length $\ell$.
    If $G$ is $\ell$-tangle-free then $B^{\ell}=B^{(\ell)}$, so
    \[
      \Pr\bigl[\,|\mu_2(B)|\text{ too large}\,\bigr]\;\le\;
      \Pr\bigl[\,G\text{ is not $\ell$-tangle-free}\,\bigr]\;+\;\Pr\bigl[\,\lVert B^{(\ell)}\rVert\text{ too large on }\one^{\perp}\,\bigr].
    \]
    \vspace{-5mm}
  \end{proposition}
  \vspace{-1mm}
  \intuition Change what is counted, not the probability space.
\end{frame}

% ------------------------------------------------------------------
\begin{frame}{Three proofs of Friedman's theorem}
  \small
  \begin{center}
  \begin{tabular}{@{}L{0.16\textwidth}L{0.25\textwidth}L{0.25\textwidth}L{0.25\textwidth}@{}}
    \toprule
    & \textbf{Friedman} (2003) & \textbf{Bordenave} (2015) & \textbf{Polynomial method} (2024)\\
    \midrule
    Object & Counts of closed non-backtracking walks & Operator norm of powers of $B$ & Expected traces of polynomials in $A$\\
    \addlinespace
    Tangles & Classified exactly by $h(\psi)$ against $\sqrt{d-1}$ & Any ball containing two cycles & Not removed; visible in higher-order terms\\
    \addlinespace
    Trivial eigenvalue & Cancelled term by term in a $1/n$ expansion & Edge indicators centred by their mean & Trace is a rational function of $1/n$\\
    \addlinespace
    Failure probability & Optimal power of $1/n$ & Polynomially small & Optimal in the permutation model\\
    \addlinespace
    Length & $\approx120$ pages & $\approx20$ pages & Short and general\\
    \bottomrule
  \end{tabular}
  \end{center}
\end{frame}

% ------------------------------------------------------------------
\begin{frame}{Summary}
  \begin{itemize}
    \setlength{\itemsep}{3mm}
    \item \textbf{Barrier.} $\edge$ is the growth rate of closed walks in $T_d$, and every $d$-regular
          graph has at least that many (Alon--Boppana).
    \item \textbf{Theorem.} Random $d$-regular graphs attain the barrier up to $\eps$ (Friedman).
    \item \textbf{Obstruction.} The trace method needs walks of length $k\gg\log n$, and at that length
          rare tangles dominate the expectation.
    \item \textbf{Resolution.} Count only tangle-free walks, and bound the probability that a tangle exists separately.
  \end{itemize}
\end{frame}

% ==================================================================
\section{Beyond the literature}

\begin{frame}{Graph lifts}
  \begin{definition}
    Let $H$ be a finite graph and $N\ge1$. An \emph{$N$-lift} of $H$ has vertex set $V(H)\times[N]$;
    for each edge $\{u,v\}$ of $H$ one chooses a permutation $\pi_{uv}$ of $[N]$ and joins $(u,i)$ to
    $(v,\pi_{uv}(i))$ for every $i$. The set $\{u\}\times[N]$ is the \emph{fibre} over $u$. In a
    \emph{random} $N$-lift the permutations are independent and uniform.
  \end{definition}
  \begin{columns}[c,onlytextwidth]
    \begin{column}{0.50\textwidth}
      \centering
      % The cube as a 2-lift of K_4: each fibre is a pair of antipodal vertices.
% Requires: \liftunit (unit length), colours ink and accent.
\begin{tikzpicture}[x=\liftunit, y=\liftunit,
    e/.style={black!55, semithick},
    f1/.style={circle, draw=ink, fill=ink, text=white, inner sep=1.4pt, font=\footnotesize},
    f2/.style={circle, draw=accent, fill=accent, text=white, inner sep=1.4pt, font=\footnotesize},
    f3/.style={circle, draw=ink, fill=white, text=ink, inner sep=1.4pt, font=\footnotesize},
    f4/.style={circle, draw=ink, fill=black!18, text=ink, inner sep=1.4pt, font=\footnotesize}]
  % base graph K_4
  \begin{scope}
    \node[f1] (a) at (0,0) {1}; \node[f2] (b) at (2,0) {2};
    \node[f4] (c) at (2,2) {4}; \node[f3] (d) at (0,2) {3};
    \draw[e] (a)--(b)--(c)--(d)--(a) (a)--(c) (b)--(d);
    \node[font=\footnotesize, anchor=north] at (1,-0.5) {base graph $K_4$};
  \end{scope}
  \draw[->, black!55] (3.0,1) -- (4.4,1);
  % the cube: vertices are binary triples, antipodal pairs form the fibres
  \begin{scope}[xshift={5.4*\liftunit}, yshift={-0.6*\liftunit}]
    \node[f1] (v000) at (0,0) {1};     \node[f2] (v100) at (3.2,0) {2};
    \node[f4] (v110) at (3.2,3.2) {4}; \node[f3] (v010) at (0,3.2) {3};
    \node[f4] (v001) at (1,1) {4};     \node[f3] (v101) at (2.2,1) {3};
    \node[f1] (v111) at (2.2,2.2) {1}; \node[f2] (v011) at (1,2.2) {2};
    \draw[e] (v000)--(v100)--(v110)--(v010)--(v000);
    \draw[e] (v001)--(v101)--(v111)--(v011)--(v001);
    \draw[e] (v000)--(v001) (v100)--(v101) (v110)--(v111) (v010)--(v011);
    \node[font=\footnotesize, anchor=north] at (1.6,-0.5) {a $2$-lift: the cube};
  \end{scope}
\end{tikzpicture}

    \end{column}
    \begin{column}{0.46\textwidth}
      \textbf{Example.} The cube is a $2$-lift of $K_4$. Each fibre (vertices with the same label)
      is a pair of antipodal vertices.

      \textbf{Remarks.} A lift of a $d$-regular graph is $d$-regular, and every vertex of the
      fibre over $u$ has exactly one neighbour in the fibre over each neighbour of $u$.
    \end{column}
  \end{columns}
\end{frame}

% ------------------------------------------------------------------

\begin{frame}{The problem: recover colors that were erased}
  \begin{columns}[c,onlytextwidth]
  \begin{column}{.45\textwidth}\centering
  \begin{tikzpicture}[scale=.88,every node/.style={circle,draw=ink,fill=white,minimum size=5mm,inner sep=0pt}]
  \foreach \v/\x/\y in {1/0/1.5,2/1.8/1.5,3/.9/3,4/0/0,5/1.8/0,6/.9/-1.5}
    \node (v\v) at (\x,\y) {};
  \foreach \u/\v in {1/2,2/3,3/1,4/5,5/6,6/4,1/4,2/5,3/6}\draw[ink] (v\u)--(v\v);
  \end{tikzpicture}\par Graph supplied
  \end{column}
  \begin{column}{.08\textwidth}\centering $\Longrightarrow$\end{column}
  \begin{column}{.45\textwidth}\centering
  \begin{tikzpicture}[scale=.88,every node/.style={circle,draw=ink,minimum size=5mm,inner sep=0pt}]
  \foreach \v/\x/\y in {1/0/1.5,2/1.8/1.5,3/.9/3}\node[fill=accent!70] (v\v) at (\x,\y) {};
  \foreach \v/\x/\y in {4/0/0,5/1.8/0,6/.9/-1.5}\node[fill=ink!50] (v\v) at (\x,\y) {};
  \foreach \u/\v in {1/2,2/3,3/1,4/5,5/6,6/4,1/4,2/5,3/6}\draw[ink] (v\u)--(v\v);
  \end{tikzpicture}\par Hidden groups recovered
  \end{column}
  \end{columns}
  \small Every vertex has \textbf{two neighbors of its own color, one of the other}.
  The drawing illustrates the rule; the theorem concerns growing random graphs.
  \textbf{Given:} the graph and fixed bounds $K,\Delta$.
  \textbf{Unknown:} colors, number of groups, sizes, and the count table.
\end{frame}

\begin{frame}{Exact counts can work when a spectral gap cannot}
  Give red vertices weight $+1$ and blue vertices weight $-1$.
  Two same-color neighbors and one other give $2-1=1$: \quad $As=s$.
  \begin{center}
  \begin{tikzpicture}[x=1.3cm,y=1cm]
    \fill[ink!12] (-2.828,-.3) rectangle (2.828,.3);
    \draw[->] (-3.4,0)--(3.6,0);
    \foreach \x/\lab in {-2.828/{-2\sqrt2},0/0,2.828/{2\sqrt2}} {\draw (\x,-.12)--(\x,.12);\node[below] at (\x,-.3) {$\lab$};}
    \fill[accent] (1,0) circle (3pt);\node[above,accent] at (1,.4) {hidden-group eigenvalue $1$};
    \node[below] at (0,-.85) {Inside the cubic spectral bulk: no isolated outlier};
  \end{tikzpicture}
  \end{center}
  \small \textbf{Model:} independent uniform simple regular/biregular factors;
  fixed positive class proportions and connected type support; exact counts.
  \textbf{Certificate:} symmetric $0\le q_{ij}<1/2$ with $\sum_jd_{ij}q_{ij}=1$.
  Every fixed regular lift template of degree $d\ge3$ qualifies ($q=1/d$).
  No noise guarantee; minimum degree three alone is insufficient.
\end{frame}

\begin{frame}{The key step: find every modular solution, then check}
  \begin{center}
  \begin{tikzpicture}[every node/.style={draw=ink,rounded corners,align=center,text width=2.25cm,minimum height=1.05cm,font=\small}]
  \node (a) at (0,0) {Try each bounded\\count table $R$};
  \node (b) at (2.9,0) {Solve $AY=YR$\\modulo a prime $p$};
  \node (c) at (5.8,0) {Enumerate the\\complete kernel};
  \node (d) at (8.7,0) {Check exact\\integer counts};
  \draw[->,thick] (a)--(b);\draw[->,thick] (b)--(c);\draw[->,thick] (c)--(d);
  \end{tikzpicture}
  \end{center}
  \textbf{A compact description replaces guessing all colorings.}
  $Y$ has one $1$ per row, recording that vertex's group.
  A sparse kernel routine certifies that \emph{no solutions were missed}.

  If its dimension is $h_R$, there are $p^{h_R}$ modular candidates.
  The essential estimate is
  \[\sup_n\E p^{h_R(A)}<\infty.\]
  \small This controls the \emph{whole expected search}, including rare large kernels.
  Integer checks reject false modular solutions. Listing is always correct.
\end{frame}

\begin{frame}{The result: exact recovery with a charged expected cost}
  \setlength{\parskip}{0.3em}
  \begin{block}{Fixed $K,\Delta$; the certified random model}
  List every valid partition in expected $\widetilde O_{K,\Delta}(n^2)$ bit work,
  expected $\widetilde O_{K,\Delta}(n)$ working bits, and expected $O(n)$ output bits.
  With high probability the \textbf{finest listed answer is the hidden partition}.
  \end{block}
  \small \textbf{Improvement in scope.} The earlier two-group spectral guarantee needs
  $(a-b)^2>4(a+b-1)$. The case $(2,1)$ fails that test; our exact-count
  theorem covers it, without knowing the table or assuming an eigengap.

  \textbf{Runtime comparison.} Sparse certified kernels avoid the $O(n^3)$
  cost of classical dense elimination. This is an expected bound in a fixed model,
  not a measured speedup or a global fastest-algorithm claim.

  \textbf{Credit and limits.} Algebraic fiber recovery has 1993/2023 precedents.
  Exact counts, fixed parameters, independent factors; no robustness to noise.
  \par\vspace{1mm}\scriptsize
  Full statements and proofs: \href{report-long.pdf}{report-long.pdf}, Part III.
  Prior spectral guarantee: Brito et al., SODA 2016 (arXiv:1507.00930).
  Algebraic precedent: Bernstein--Cottaar--Lange, IPCC report (2023), Section 3.1.
\end{frame}

% ------------------------------------------------------------------
\begin{frame}{References}
  \scriptsize
  \begin{thebibliography}{9}
    \setbeamertemplate{bibliography item}[text]
    \setlength{\itemsep}{0.5mm}
    \bibitem{A} N.~Alon. \emph{Eigenvalues and expanders}. Combinatorica 6 (1986).
    \bibitem{N} A.~Nilli. \emph{On the second eigenvalue of a graph}. Discrete Math.\ 91 (1991).
    \bibitem{M} B.~D.~McKay. \emph{The expected eigenvalue distribution of a large regular graph}.
      Linear Algebra Appl.\ 40 (1981).
    \bibitem{F} J.~Friedman. \emph{A proof of Alon's second eigenvalue conjecture and related problems}.
      Mem.\ Amer.\ Math.\ Soc.\ 195 (2008). arXiv:cs/0405020.
    \bibitem{FK} J.~Friedman and D.-E.~Kohler. \emph{The relativized second eigenvalue conjecture of Alon}.
      arXiv:1403.3462 (2014).
    \bibitem{B} C.~Bordenave. \emph{A new proof of Friedman's second eigenvalue theorem and its extension to random lifts}.
      Ann.\ Sci.\ \'Ec.\ Norm.\ Sup\'er.\ 53 (2020). arXiv:1502.04482.
    \bibitem{BDH} G.~Brito, I.~Dumitriu and K.~D.~Harris. \emph{Spectral gap in random bipartite biregular graphs and applications}.
      Combin.\ Probab.\ Comput.\ 31 (2022). arXiv:1804.07808.
    \bibitem{CGTV} C.-F.~Chen, J.~Garza-Vargas, J.~A.~Tropp and R.~van Handel. \emph{A new approach to strong convergence}.
      arXiv:2405.16026 (2024).
  \end{thebibliography}
\end{frame}

\end{document}
